% ============================================================================
% Fichier  : front/01_avant_propos.tex
% Titre    : Avant-propos
% Auteur   : MINKA MI NGUIDJOÏ Thierry Emmanuel
% Version  : 1.0 -- novembre 2026
% ============================================================================
% Décision du 06/11/2026 : seuls A0 et A1 (HAL) sont cités ; aucun travail
% non publié n'est annoncé. 60 h = 30 séances de 2 h ; 31 chapitres.
% ============================================================================

\chapter*{Avant-propos}
\addcontentsline{toc}{chapter}{Avant-propos}
\markboth{Avant-propos}{Avant-propos}

\epigraph{« Ce que l'organisation ignore de ses propres systèmes conditionne
la qualité des décisions qu'elle peut prendre à leur sujet. »}%
{\textit{Vers l'auditologie}~\cite{minka2026auditologie}}

Pour la plupart des élèves ingénieurs du Département de Génie Informatique
de l'École Nationale Supérieure Polytechnique de Yaoundé, ce cours est le
seul consacré à l'audit, au contrôle interne, à la gestion des risques, à
l'audit de sécurité et à l'audit de conformité. Il est donc
autoportant~: il ne suppose ni cours de gouvernance, ni cours de gestion
des risques préalable. Il prépare aussi aux UE d'investigation numérique
(SEC4031 et SEC4052), qui reprennent là où il s'arrête~: à la preuve.

Il repose sur une double conviction. La première est qu'auditer ne consiste
pas à remplir une grille, mais à réduire l'incertitude d'un tiers sur un
système qu'il ne peut pas observer seul. Un conseil d'administration, un
régulateur, un juge s'appuient sur ce que l'auditeur affirme~; ils doivent
pouvoir savoir aussi ce qu'il ignore. C'est pourquoi ce cours admet, à côté
de \OUI{} et de \NON{}, une troisième réponse~: \NEK{}, « connaissance
insuffisante ».

La seconde conviction est que le métier ne s'importe pas tel quel. Les
normes de l'IIA, de l'ISACA, de l'ISO et de l'INTOSAI sont des acquis
précieux~; elles s'appliquent ici sous le regard de la loi camerounaise sur
la cybersécurité et la cybercriminalité, de la COBAC, du droit OHADA et de
la Chambre des comptes, dans des organisations dont les moyens et les
contraintes sont les nôtres.

\bigskip

\textbf{Ce que ce cours n'est pas.} Il n'est pas un résumé des normes
professionnelles, ni une préparation à une certification, ni un cours de
tests d'intrusion. Les normes y sont citées avec exactitude, dans l'édition
en vigueur quand elle a pu être vérifiée~; les textes de loi n'y sont cités
entre guillemets que s'ils ont été lus dans leur version officielle.

\bigskip

\textbf{Ce que ce cours est.} Un traité de trente et un chapitres, en six
parties précédées d'un prologue, qui va du cas concret vers la méthode, de
la méthode vers la pratique, et de la pratique vers la recherche. La
Partie~I pose les fondements~: audit, contrôle interne, risque et sa
gestion, fraude. La Partie~II donne la méthode de la mission, de la lettre
de mission au suivi. La Partie~III présente les normes, les référentiels,
le droit et la déontologie. La Partie~IV audite le système d'information
domaine par domaine, de la gouvernance à la conformité, en passant par la
sécurité, les tests d'intrusion, la continuité et les tiers, et se termine
par trois cas intégrateurs. La Partie~V aborde l'audit continu,
l'intelligence artificielle et le passage de l'audit à la preuve
judiciaire. La Partie~VI présente le programme de recherche sur lequel la
méthode s'appuie.

\bigskip

\textbf{Un cours, un corpus.} La méthode enseignée est l'instanciation
opérationnelle, pour le praticien, du programme de recherche de
l'auditologie~\cite{minka2026auditologie}~; son instrument de réponse
\OUI{}~/ \NON{}~/ \NEK{} et ses mesures viennent de l'article
\cite{minka2026nek}. Ces deux textes sont des prépublications déposées sur
HAL, et le cours les cite comme telles. Les autres appuis sont des travaux
publiés et vérifiables~: théorie de l'information, théorie des jeux,
psychologie cognitive, métrologie, normes professionnelles. Ce cours est la
version complète du corpus~; le cours ACSI du Master de la Chaire UNESCO en
est une version resserrée.

\bigskip

\textbf{Trois organisations fictives.} Le cours suit trois missions.
\emph{LiZbiZ}, entreprise de logistique aux quatre centres de données,
reprise de l'article fondateur, est le cas principal~: le cabinet AUDIX
CONSEIL y conduit deux cycles d'audit. La \emph{BANQUE SANAGA}, à Douala,
affronte les suites d'une panne et le regard de la COBAC. L'\emph{ANPP},
agence publique de paiement, est auditée par la Chambre des comptes. Les
noms, les dates et les montants sont fictifs~; les situations sont
plausibles, et chaque chiffre se recalcule.

\bigskip

\textbf{Comment travailler.} Lisez chaque chapitre avec sa fiche
d'auto-évaluation~; refaites les calculs des lectures auditologiques~;
faites les labs de NDAP SUKULU avant d'en lire les corrigés. Surtout,
prenez l'habitude que le cours cherche à transmettre~: séparer ce que l'on
a constaté de ce que l'on en pense, et dire ce que l'on ne sait pas.

\bigskip

\begin{flushright}
    \textit{\auteurcours}\\
    \textit{Yaoundé, novembre 2026}
\end{flushright}
